\section{Estado del Arte}
\label{sec:estado-arte}

Esta sección revisa los marcos teóricos y las herramientas en las que
se apoya el sistema, distinguiendo aquellas que han sido finalmente
adoptadas de las alternativas evaluadas y descartadas.

\subsection{Orquestación de agentes con LangGraph}

LangGraph~\cite{langgraph} (LangChain Labs) es una biblioteca para
construir sistemas multiagente como grafos dirigidos cuyos nodos son
funciones Python y cuyas aristas pueden ser condicionales,
ejecutadas en función del estado tipado del grafo. Frente a
alternativas como \texttt{AutoGen} (Microsoft) o \texttt{CrewAI},
LangGraph destaca por:

\begin{itemize}[leftmargin=1.5em,itemsep=0.2em]
  \item Su \textbf{enfoque grafo + estado tipado}, que permite
        razonar formalmente sobre la trayectoria del control y
        garantiza la \textit{terminación} (en este proyecto se
        impone un \textit{iteration limit} de seis pasos).
  \item Su \textbf{soporte nativo de checkpointers} (\texttt{MemorySaver},
        \texttt{PostgresSaver}, \texttt{SqliteSaver}), que guarda
        el histórico conversacional dentro del mismo grafo sin
        necesidad de un sistema externo.
  \item La posibilidad de combinar \textit{structured output}
        (Pydantic) con LLM \textit{plain} en distintos nodos del
        mismo grafo. En el sistema se aprovecha para separar
        \textit{routing} (decisión) y \textit{narración} (texto
        natural).
\end{itemize}

\subsection{LLM y \textit{providers}}

El catálogo de proveedores ha quedado configurable por usuario, con
soporte nativo para Groq~\cite{groq}, OpenAI, Anthropic y Google. El
modelo por defecto es \texttt{llama-3.3-70b-versatile} servido por
Groq, que ofrece latencias inferiores a 500\,ms en \textit{tool
calling} gracias a su hardware LPU dedicado. Para tareas de
\textit{small-talk} y resumen rodante se usa
\texttt{llama-3.1-8b-instant}, que reduce el gasto de tokens en
operaciones poco exigentes.

\subsection{Visión por computador: OCR y biometría}

Para OCR se adoptó \texttt{PaddleOCR}~\cite{paddleocr} (versión 2.6.2)
por su rendimiento sobre texto en castellano y su capacidad de
funcionar en CPU sin GPU externa, descartando \texttt{EasyOCR} (menor
precisión sobre tickets en euros). Para biometría se eligió
\texttt{facenet-pytorch}~\cite{facenet} (\texttt{InceptionResnetV1}
pre-entrenado en VGGFace2) en lugar de \texttt{deepface}/\texttt{ArcFace},
evitando así introducir TensorFlow en el stack. La detección de rostros se apoya en
\texttt{MTCNN}~\cite{mtcnn} y la detección de vida (\textit{liveness})
en \texttt{DenseNet201}~\cite{densenet} con pesos ImageNet, completada
con un módulo \textit{anti-spoofing} basado en LBP, análisis de espacio
de color HSV/YCrCb y análisis frecuencial (FFT) para detectar
patrones Moiré de pantallas LCD. Los \textit{landmarks} faciales
necesarios para el \textit{challenge-response} de parpadeo se extraen
con \texttt{MediaPipe Face Mesh}~\cite{mediapipe}.

\subsection{Detección de anomalías y predicción temporal}

El detector de anomalías financieras combina
\texttt{IsolationForest}~\cite{isolation_forest} de
\texttt{scikit-learn} con una regla heurística de 3-sigma sobre los
importes históricos del usuario. Esta combinación, ya validada en P5,
se ha mantenido por su robustez y por la interpretabilidad de la
regla 3-sigma frente al usuario final. Para la predicción mensual
conviven tres modelos --- Random Forest, HistGradientBoosting y ARIMA
--- accesibles vía un único endpoint REST que permite seleccionar el
método.

\subsection{Clasificación de texto financiero}

Para clasificar la \texttt{Area} de una transacción se ha evolucionado
del \texttt{SGDClassifier} sobre TF-IDF de char $n$-gramas (P2
original) hacia un pipeline híbrido basado en
\textit{Sentence-BERT}~\cite{sentence_bert} con embeddings
multilingües. Este enfoque permite generalizar a descripciones que
no aparecieron en el entrenamiento (\textit{cold-start} y nuevos
idiomas) sin perder el \textit{macro-F1} en el conjunto de
entrenamiento.

\subsection{Observabilidad: Langfuse}

\texttt{Langfuse}~\cite{langfuse} (versión 3) es una plataforma de
observabilidad específica para LLM con instrumentación nativa de
LangChain/LangGraph mediante \texttt{CallbackHandler}. Recoge
automáticamente la cadena de llamadas (\textit{traces}), prompt
completo, respuesta cruda, \textit{tool calls}, latencias y consumo
de tokens. Frente a alternativas como \texttt{LangSmith} (de pago),
\texttt{Helicone} (no LangChain-native) o \texttt{Phoenix} (Arize),
Langfuse ofrece un plan gratuito generoso y un \textit{dashboard} web
suficiente para auditar la ejecución de los agentes en tiempo real.

\subsection{Prompt engineering: ASPECCT vs TAREA}

El enunciado permite escoger entre dos metodologías estructuradas de
\textit{prompting}. ASPECCT~\cite{aspect} (Audience, Style, Purpose,
Examples, Constraints, Context, Tone) separa \textit{style} y
\textit{tone}, lo que en el sistema permite inyectar el bloque de
estilo dependiente del rol del usuario
(\texttt{basic}/\texttt{advanced}) como un \texttt{SystemMessage}
independiente sin tocar el prompt principal. TAREA (Tarea, Audiencia,
Recursos, Estructura, Activación) resulta más compacta pero no separa
estilo y tono. Por esta razón se ha adoptado ASPECCT como metodología
oficial del proyecto.

\subsection{Pruebas de carga: Locust}

Para validar el comportamiento del sistema bajo concurrencia se ha
optado por \texttt{Locust}~\cite{locust}, un \textit{framework} de
\textit{load testing} en Python que permite describir el
comportamiento de los usuarios virtuales como clases con \texttt{@task}.
La principal ventaja sobre alternativas como JMeter o k6 es que los
escenarios se escriben en el mismo lenguaje del backend, lo que
facilita la generación de cargas realistas (registro con
\textit{multipart}, JWT, etc.) sin duplicar lógica.
